%<*driver>
\def\nameofplainTeX{plain}
\ifx\fmtname\nameofplainTeX
\input l3docstrip.tex
\keepsilent
\askforoverwritefalse
\preamble

========================================================================
 Package tagbridge - Tagging-compatibility hooks for packages without
 native PDF/UA support.

 (C) 2026 Matthias Werner
 E-mail: matthias.werner@informatik.tu-chemnitz.de

 Released under the LaTeX Project Public License v1.3c or later
 See http://www.latex-project.org/lppl.txt
========================================================================

\endpreamble
\postamble

 This work is "maintained" (as per LPPL maintenance status) by
 Matthias Werner.
\endpostamble
\generate{\file{tagbridge.sty}{\from{tagbridge.dtx}{package}}}
\expandafter\endbatchfile
\fi

\iffalse
%</driver>
%<*package>
\NeedsTeXFormat{LaTeX2e}[2022-06-01]
\ProvidesExplPackage{tagbridge}{2026-09-14}{v0.1.0}
  {Tagging-compatibility hooks for packages without native PDF/UA support}
%</package>
%<*driver>
\fi
\documentclass[add-bib=false]{osgdoc}

\OnlyDescription
\usepackage{enumitem}
\usepackage{microtype}
\usepackage[
languages={de,en},
map={de=german,en=english},
targetlang={job=2},
load babel
]{langselect}
\newpackagename{\tagbridgepkg}{tagbridge}
\makeindex
\setcnltx{
  package  =  tagbridge,
  version  = \UseVersionOf{tagbridge.sty},
  date     = \UseDateOf{tagbridge.sty},
  title    = {\ldeen{Das \tagbridgepkg\ Paket}{The \tagbridgepkg\ Package}},
  info     = {\ldeen{Tagging-Kompatibilitätshooks für Fremdpakete ohne
      eigene PDF/UA-Unterstützung}{Tagging-compatibility hooks for
      third-party packages without native PDF/UA support}},
  authors  = {Matthias Werner[matthias.werner@informatik.tu-chemnitz.de]},
  abstract={
    \ldeen*{
      \tagbridgepkg\ sammelt an einem Ort, was einzelne Fremdpakete
      brauchen, um unter \cs{DocumentMetadata}\Marg{tagging=on} sauber
      getaggten statt fehlerhaften oder anonymen PDF-Output zu erzeugen.
      Jeder Hook greift nur, wenn \pkg{tagpdf} aktiv ist (getestet über
      \cs{tagstructbegin}), und bleibt ansonsten wirkungslos. Aktuell
      enthalten: ein Hook für \pkg{qrcode} (das seine Module einzeln über
      \cs{rule}s zeichnet und ohne Eingriff Hunderte Tagging-Warnungen pro
      Code erzeugt) sowie das Nachladen von
      \pkg*{latex-lab-testphase-tikz} (das \env{tikzpicture} die Schlüssel
      \code{alt}/\code{artifact} gibt, die \pkg{pgf}/\pkg{tikz} selbst nur
      reserviert, aber nie auswertet).

      Das Paket ist als Übergangslösung gedacht: sobald das
      \pkg*{LaTeX-Project} selbst -- über \pkg*{latex-lab} oder dessen
      Nachfolger -- native Unterstützung für ein hier behandeltes Paket
      liefert, wird der entsprechende Hook hier überflüssig und kann
      entfernt werden. \pkg*{pdfpages}-Import fällt bewusst nicht in
      diesen Zuständigkeitsbereich; dafür existiert bereits
      \pkg*{tagpax}.
    }{
      \tagbridgepkg\ collects in one place what individual third-party
      packages need in order to produce cleanly tagged rather than broken
      or anonymous PDF output under \cs{DocumentMetadata}\Marg{tagging=on}.
      Each hook only fires while \pkg{tagpdf} is active (tested via
      \cs{tagstructbegin}) and stays inert otherwise. Currently included:
      a hook for \pkg{qrcode} (which draws its modules individually via
      \cs{rule}s and, left alone, produces hundreds of tagging warnings
      per code) and loading \pkg*{latex-lab-testphase-tikz} (which gives
      \env{tikzpicture} the \code{alt}/\code{artifact} keys that
      \pkg{pgf}/\pkg{tikz} itself only reserves but never evaluates).

      The package is meant as a stopgap: once the \pkg*{LaTeX-Project}
      itself -- via \pkg*{latex-lab} or its successor -- ships native
      support for a package handled here, the corresponding hook becomes
      redundant and can be removed. Importing PDFs via \pkg*{pdfpages} is
      deliberately out of scope here; \pkg*{tagpax} already covers that.
    }
  },
  url      =https://github.com/werner-matthias/osglecture,
  build-title
}
\begin{document}

\section{\ldeen{Zweck}{Purpose}}

\ldeen{Das Paket wird üblicherweise nicht direkt geladen, sondern von
einer Dokumentklasse oder einem anderen Paket (zum Beispiel
\pkg*{osglecture}) als Abhängigkeit eingebunden. Es hat keine
Paketoptionen und keine öffentlichen Befehle -- alle Hooks greifen
automatisch, sobald das jeweils betroffene Fremdpaket
(\pkg{qrcode}, \pkg{tikz}) geladen wird, und nur, wenn
\pkg{tagpdf}-Tagging aktiv ist.}{The package is usually not loaded
directly but pulled in as a dependency by a document class or another
package (for example \pkg*{osglecture}). It has no package options and
no public commands -- every hook fires automatically once the
respective third-party package (\pkg{qrcode}, \pkg{tikz}) loads, and
only while \pkg{tagpdf} tagging is active.}

\ldeen{Die \pkg{qrcode}-Optionen \code{tag} (Default \code{Figure}) und
\code{alt} (Default: sprachabhängiger Text mit dem codierten Inhalt)
lassen sich wie gewohnt pro Aufruf oder via \cs{qrset} setzen;
\code{tag=artifact} markiert den Code als rein dekoratives Artifact.}
{The \pkg{qrcode} options \code{tag} (default \code{Figure}) and
\code{alt} (default: language-dependent text carrying the encoded
content) can be set per call or via \cs{qrset} as usual;
\code{tag=artifact} marks the code as a purely decorative artifact.}

\section{Implementation}
\MaybeStop{\ldeen*{Um die Implementationsbeschreibung zu erhalten, kommentieren Sie bitte \cs{OnlyDescription} in @1 aus.}{To view the implementation description, please uncomment \cs{OnlyDescription} in @1.}{\File{tagbridge.dtx}}}

\AutoDocInput{tagbridge.dtx}{package}
\end{document}

%</driver>

%<*package>
\ExplSyntaxOn

% \ldeen*{@1 zeichnet jedes Modul des Codes einzeln über \cs{rule}s statt
% über ein Bild-XObject. Bei aktivem \pkg{tagpdf} greift dafür in der
% Lua-Variante der automatische \code{activate/tagunmarked}-Mechanismus
% und wrappt jedes sonst ungetaggte Modul separat als Artifact -- mit
% einer eigenen Warnung pro Modul, bei einem mittelgroßen Code schnell
% mehrere hundert. Der Hook patcht \cs{qr@printmatrix} und
% \cs{qr@printsavedbinarymatrix} (die einzigen beiden Stellen, an denen
% @1 tatsächlich zeichnet), sodass der komplette Code in \emph{einer}
% eigenen Tag-Struktur steckt, bevor der automatische Mechanismus je
% einzelne Module sieht. Der Hook hängt an \code{package/qrcode/after}
% und greift daher unabhängig davon, von wem @1 geladen wird. Die neuen
% @1-Optionen \code{tag} (Default \code{Figure}) und \code{alt} (Default:
% der codierte Text) erlauben, das pro Aufruf oder via \cs{qrset} zu
% steuern; \code{tag=artifact} markiert den Code als rein dekoratives
% Artifact ohne eigene Struktur -- sinnvoll, wenn eine Beschriftung
% daneben den Inhalt bereits als Fließtext trägt.}{@1 draws every module
% of the code individually via \cs{rule}s instead of an image XObject.
% With \pkg{tagpdf} active, the Lua variant's automatic
% \code{activate/tagunmarked} mechanism steps in and wraps each
% otherwise-untagged module separately as an artifact -- with its own
% warning per module, quickly reaching several hundred for a
% medium-sized code. The hook patches \cs{qr@printmatrix} and
% \cs{qr@printsavedbinarymatrix} (the only two places where @1 actually
% draws) so that the whole code sits inside \emph{one} tag structure
% before the automatic mechanism ever sees an individual module. The
% hook attaches to \code{package/qrcode/after} and therefore fires
% regardless of who loads @1. The new @1 options \code{tag} (default
% \code{Figure}) and \code{alt} (default: the encoded text) let this be
% controlled per call or via \cs{qrset}; \code{tag=artifact} marks the
% code as a purely decorative artifact with no structure of its own --
% useful when a caption next to it already carries the content as flow
% text.}{\pkg{qrcode}}
% \ldeen*{\code{alt}-Default und Kritikprüfung: Ohne explizites
% \code{alt} übernimmt \cs{__tagbridge_qrcode_alt_text:} den codierten
% Text sprachabhängig -- über \cs{olsTargetLanguage}, falls das existiert
% (z.\,B.\ von \pkg*{langselect} gesetzt), sonst über
% \cs{OsgLectureLanguage}, falls \pkg*{osglecture} geladen ist, sonst
% \code{en}. Beide Integrationen sind rein optional: ohne sie fällt der
% Default einfach auf Englisch zurück, das Paket funktioniert auch ohne
% \pkg*{osglecture} oder \pkg*{langselect}. Für Inhalte, die mit einem
% der wenigen tatsächlich \emph{geheimnistragenden} Protokollpräfixe
% beginnen (\code{WIFI:}, \code{otpauth://} -- WLAN-Passwort bzw.\
% TOTP-Seed), würde dieser Default das Geheimnis im Klartext in den
% Tag-Baum des PDFs schreiben. Andere strukturierte Inhalte wie
% \code{MECARD:}/\code{BEGIN:VCARD} tragen zwar persönliche, aber keine
% Zugangsdaten und bleiben bewusst außen vor. Für die kritische
% Teilmenge bricht \cs{__tagbridge_qrcode_check_critical_content:} mit
% einem Fehler ab, statt den Default stillschweigend zu verwenden --
% explizites \code{alt=\{...\}} (das den Inhalt beschreibt, ohne ihn zu
% wiederholen) oder \code{tag=artifact} ist dann erforderlich.}{Default
% for \code{alt} and criticality check: without an explicit \code{alt},
% \cs{__tagbridge_qrcode_alt_text:} falls back to the encoded text,
% language-dependent -- via \cs{olsTargetLanguage} if that exists (e.g.\
% set by \pkg*{langselect}), otherwise via \cs{OsgLectureLanguage} if
% \pkg*{osglecture} is loaded, otherwise \code{en}. Both integrations are
% purely optional: without them the default simply falls back to
% English, the package works fine without \pkg*{osglecture} or
% \pkg*{langselect}. For content starting with one of the few protocol
% prefixes that actually carry a secret (\code{WIFI:}, \code{otpauth://}
% -- Wi-Fi password resp.\ TOTP seed), this default would write that
% secret in the clear into the PDF's tag tree. Other structured content
% such as \code{MECARD:}/\code{BEGIN:VCARD} carries personal but not
% credential data and is deliberately left out. For the critical subset,
% \cs{__tagbridge_qrcode_check_critical_content:} aborts with an error
% instead of silently using the default -- an explicit \code{alt=\{...\}}
% (describing the content without repeating it) or \code{tag=artifact}
% is then required.}
\cs_if_exist:NT \tagstructbegin
  {
    \tl_new:N \l__tagbridge_qrcode_tag_tl
    \tl_new:N \l__tagbridge_qrcode_alt_tl
    \tl_new:N \l__tagbridge_qrcode_tmpa_tl
    \msg_new:nnn { tagbridge } { qrcode-critical-content }
      {
        This~QR~code's~content~starts~with~WIFI:~or~otpauth://~--~a~
        format~that~typically~carries~a~credential~or~secret~(Wi-Fi~
        password,~one-time-password~seed,~...).~Falling~back~to~the~
        encoded~text~as~the~PDF~alt~text~would~leak~it~into~the~
        document's~accessibility~tree,~where~it~is~trivially~
        extractable.~Set~the~qrcode~option~alt=\{...\}~explicitly~
        (describing~the~code~without~repeating~the~secret),~or~
        tag=artifact~if~nothing~about~it~should~be~exposed~at~all.
      }
    \cs_new_protected:Npn \__tagbridge_qrcode_check_critical_content:
      {
        \tl_if_empty:NT \l__tagbridge_qrcode_alt_tl
          {
            \tl_if_eq:NnF \l__tagbridge_qrcode_tag_tl { artifact }
              {
                \tl_set:Ne \l__tagbridge_qrcode_tmpa_tl { \qr@texttoencode }
                \regex_match:nVT { \A(WIFI:|otpauth://) }
                  \l__tagbridge_qrcode_tmpa_tl
                  { \msg_error:nn { tagbridge } { qrcode-critical-content } }
              }
          }
      }
    \cs_new:Npn \__tagbridge_qrcode_alt_text:
      {
        \tl_if_empty:NTF \l__tagbridge_qrcode_alt_tl
          {
            \str_if_eq:eeTF
              {
                \cs_if_exist:NTF \olsTargetLanguage
                  { \olsTargetLanguage }
                  {
                    \cs_if_exist:NTF \OsgLectureLanguage
                      { \OsgLectureLanguage }
                      { en }
                  }
              }
              { de }
              { QR-Code:~\qr@texttoencode }
              { QR~code:~\qr@texttoencode }
          }
          { \l__tagbridge_qrcode_alt_tl }
      }
    \cs_new_protected:Npn \__tagbridge_qrcode_tag_begin:
      {
        \tagpdfparaOff
        \__tagbridge_qrcode_check_critical_content:
        \tl_if_eq:NnTF \l__tagbridge_qrcode_tag_tl { artifact }
          { \tagmcbegin{artifact} }
          {
            \exp_args:Ne \tagstructbegin
              {
                tag=\l__tagbridge_qrcode_tag_tl,
                alt={\__tagbridge_qrcode_alt_text:}
              }
            \exp_args:Ne \tagmcbegin { tag=\l__tagbridge_qrcode_tag_tl }
          }
      }
    \cs_new_protected:Npn \__tagbridge_qrcode_tag_end:
      {
        \tl_if_eq:NnTF \l__tagbridge_qrcode_tag_tl { artifact }
          { \tagmcend }
          { \tagmcend \tagstructend }
      }
    \AddToHook { package/qrcode/after } [ tagbridge/qrcode-tagging ]
      {
        \define@key{qr}{tag}{\tl_set:Nn \l__tagbridge_qrcode_tag_tl {#1}}
        \define@key{qr}{alt}{\tl_set:Nn \l__tagbridge_qrcode_alt_tl {#1}}
        \qrset{tag=Figure}
        \let\osg@qrcode@printmatrix@orig\qr@printmatrix
        \let\osg@qrcode@printsavedmatrix@orig\qr@printsavedbinarymatrix
        \def\qr@printmatrix#1{%
          \__tagbridge_qrcode_tag_begin:
          \osg@qrcode@printmatrix@orig{#1}%
          \__tagbridge_qrcode_tag_end:
        }%
        \def\qr@printsavedbinarymatrix#1{%
          \__tagbridge_qrcode_tag_begin:
          \osg@qrcode@printsavedmatrix@orig{#1}%
          \__tagbridge_qrcode_tag_end:
        }%
      }
  }

% \ldeen*{pgf/tikz reserviert die Schlüssel \code{alt}, \code{artifact}
% und \code{actualtext} für \env{tikzpicture} ausdrücklich für
% zukünftige Tagging-Unterstützung, wertet sie selbst aber nirgends
% aus -- ein damit angegebener Alt-Text verschwindet kommentarlos, und
% das ganze Bild landet als anonymes Artifact im Tag-Baum, ohne dass
% \pkg{tagpdf} das meldet. Das \pkg{latex-lab}-Bundle bringt mit
% \pkg{latex-lab-testphase-tikz} eine vom LaTeX-Projekt selbst
% gepflegte Implementierung dieser Schlüssel über das
% Tagging-Socket-System mit, die auch mehrzeiligen Node-Text korrekt
% als echten Text taggt statt als Teil der Grafik; das Laden dieses
% einen Pakets ist robuster als eine eigene Nachbildung. Die Schlüssel
% gelten für die \env{tikzpicture}-Umgebung selbst
% (\code{\textbackslash begin\{tikzpicture\}[alt=\{...\}]}), nicht für
% einzelne \cs{node}s darin -- ein \code{alt} an einem einzelnen Knoten
% bleibt wirkungslos.}{pgf/tikz reserves the keys \code{alt},
% \code{artifact} and \code{actualtext} for \env{tikzpicture}
% explicitly for future tagging support but never evaluates them
% itself -- an \code{alt} text given this way vanishes without a
% trace, and the whole picture ends up as an anonymous artifact in the
% tag tree, with no diagnostic from \pkg{tagpdf}. The \pkg{latex-lab}
% bundle ships \pkg{latex-lab-testphase-tikz}, a
% LaTeX-Project-maintained implementation of these keys via the
% tagging socket system that also tags multi-line node text correctly
% as real text rather than as part of the graphic; loading that one
% package is more robust than reimplementing it. The keys apply to the
% \env{tikzpicture} environment itself
% (\code{\textbackslash begin\{tikzpicture\}[alt=\{...\}]}), not to
% individual \cs{node}s inside it -- an \code{alt} on a single node
% remains ineffective.}{tikzpicture}
\cs_if_exist:NT \tagstructbegin
  { \RequirePackage{latex-lab-testphase-tikz} }

\ExplSyntaxOff
%</package>
